\subsection{Jump buffer}

\textit{Jump buffer} funguje podobně jako \textit{Coyote time}. Rozdíl je ten, že \textit{Jump buffer} se zaměřuje na skok při dopadu na zem. Když hra zaznamená input od hráče v momentě, kdy hráč teprve dopadá z výšky na zem, ale ještě na ní není, avšak je pár framů nad ní, tak se tento input opět zaregistruje jako skok a herní objekt se posune na vertikální ose.